\documentclass[11pt,a4paper]{beamer}
\usetheme{Madrid}
\usepackage[utf8]{inputenc}
\usepackage{amsmath}
\usepackage{amsfonts}
\usepackage{amssymb}
\usepackage{graphicx}
\usepackage{tikz}
\usepackage{booktabs}
\usepackage{array}
\usepackage{colortbl}
\usepackage{multicol}

\usefonttheme{professionalfonts}

\setbeamerfont{normal text}{size=\tiny}
\setbeamerfont{itemize/enumerate body}{size=\tiny}
\setbeamerfont{itemize/enumerate subbody}{size=\tiny}
\setbeamerfont{block body}{size=\tiny}
\setbeamerfont{block title}{size=\scriptsize}
\setbeamerfont{frametitle}{size=\small}
\setbeamerfont{title}{size=\small}
\setbeamerfont{subtitle}{size=\footnotesize}
\setbeamerfont{author}{size=\tiny}
\setbeamerfont{date}{size=\tiny}

\setlength{\leftmargini}{0.3em}
\setlength{\leftmarginii}{0.3em}
\setlength{\labelsep}{0.1em}
\setlength{\itemsep}{0.02em}
\setlength{\parskip}{0.02em}

\title[CAMS Exam Preparation]{CAMS Exam Preparation}
\subtitle{Money Laundering Risks in Nonbank Financial Services}
\author{Money Transmitters | Fintech | Insurance | Securities | Lending | Crypto}
\date{}

\setbeamercolor{title}{fg=white,bg=blue!70!black}
\setbeamercolor{frametitle}{fg=white,bg=blue!70!black}
\setbeamertemplate{navigation symbols}{}

\begin{document}

% TITLE SLIDE
\begin{frame}
\titlepage
\centering
\tiny 15 Topics | Nonbank Financial ML Risks | How to Answer
\end{frame}

% ============================================================
% SLIDE 1 - NONBANK FINANCIAL INSTITUTIONS OVERVIEW
% ============================================================
\begin{frame}{1. Nonbank Financial Institutions Overview}
\tiny
\noindent\textbf{Introduction to the topic:} Nonbank financial institutions (NBFIs) play an increasingly important role in the financial system. They include money transmitters, payment processors, fintech companies, money service businesses (MSBs), insurance companies, securities firms, investment advisers, consumer lenders, and virtual asset service providers. Unlike banks, NBFIs may not take deposits but they handle large volumes of financial transactions and client funds. Under FATF Recommendations 22 and 23, NBFIs are subject to AML/CFT obligations similar to banks. The key difference is that NBFIs may have less regulatory oversight or less mature AML programs. Criminals exploit NBFIs because they may have weaker controls, faster onboarding, or less scrutiny of cross-border transactions.

\vspace{0.1cm}
\noindent\textbf{Type of violation or warning signs:} A money transmitter has no AML program and has never filed a SAR. A payment processor allows customers to send funds to high-risk jurisdictions without due diligence. A fintech company onboards customers with minimal identity verification (just an email address). An insurance company fails to verify the source of funds for large premium payments. A securities firm allows trades through omnibus accounts without customer identification.

\vspace{0.1cm}
\noindent\textbf{Correct answer approach:} The correct answer will state that NBFIs have AML obligations including registration, CDD, transaction monitoring, and SAR filing, similar to banks. Eliminate answers that suggest NBFIs are exempt from AML rules or have reduced obligations. Remember: NBFIs are regulated under the same AML framework as banks in most jurisdictions.
\end{frame}

% ============================================================
% SLIDE 2 - MONEY TRANSMITTERS AND MSBs
% ============================================================
\begin{frame}{2. Money Transmitters and Money Service Businesses (MSBs)}
\tiny
\noindent\textbf{Introduction to the topic:} Money transmitters (also called money service businesses or MSBs) provide money transfer services domestically and internationally. They are the primary formal alternative to banks for cross-border remittances. Under the Bank Secrecy Act and FATF standards, MSBs must register with regulators, maintain AML programs, conduct CDD on customers, file CTRs for cash transactions over \$10,000, and file SARs. MSBs are particularly vulnerable to money laundering because they handle cash, process high volumes of small transactions, and serve clients in high-risk jurisdictions.

\vspace{0.1cm}
\noindent\textbf{Type of violation or warning signs:} An MSB has not registered with FinCEN or the equivalent regulator. An MSB accepts cash payments without verifying customer identity. An MSB processes transactions just below reporting thresholds (structuring). An MSB sends funds to jurisdictions with weak AML controls. An MSB has no AML program, no compliance officer, and no employee training. An MSB has never filed a SAR despite processing high volumes.

\vspace{0.1cm}
\noindent\textbf{Correct answer approach:} The correct answer will state that MSBs must register, maintain AML programs, conduct CDD, and file CTRs and SARs. Eliminate answers that suggest MSBs are exempt from registration or reporting. Remember: MSBs are the primary gatekeepers for international money transfers — they must know their customers and the purpose of transfers.
\end{frame}

% ============================================================
% SLIDE 3 - PAYMENT PROCESSORS AND MERCHANT ACQUIRERS
% ============================================================
\begin{frame}{3. Payment Processors and Merchant Acquirers}
\tiny
\noindent\textbf{Introduction to the topic:} Payment processors and merchant acquirers facilitate electronic payments between merchants and customers. They handle credit card, debit card, ACH, and digital wallet transactions. Payment processors are vulnerable because they may not have direct visibility into the underlying merchants they serve. Criminals set up fake merchant accounts, process payments for illegal goods, or use payment processors to layer funds. A "high-risk" merchant (e.g., adult entertainment, gaming, or crypto) may be more likely to launder money. Payment processors must conduct due diligence on their merchants, monitor transactions, and file SARs for suspicious activity.

\vspace{0.1cm}
\noindent\textbf{Type of violation or warning signs:} A payment processor onboards a high-risk merchant without due diligence. A merchant receives an unusually high volume of chargebacks (indicating fraud). A merchant processes transactions that are significantly above the average transaction amount for that industry. A merchant processes payments for a business that has no website or physical address. A payment processor does not monitor merchants for suspicious activity. A payment processor has never filed a SAR.

\vspace{0.1cm}
\noindent\textbf{Correct answer approach:} The correct answer will state that payment processors must conduct due diligence on merchants, monitor transactions, and file SARs. Eliminate answers that suggest payment processors have no responsibility for merchant behavior. Remember: payment processors are gatekeepers for electronic payments — they must know who they are serving.
\end{frame}

% ============================================================
% SLIDE 4 - FINTECH AND DIGITAL PAYMENT PLATFORMS
% ============================================================
\begin{frame}{4. Fintech and Digital Payment Platforms}
\tiny
\noindent\textbf{Introduction to the topic:} Fintech companies offer innovative financial services including mobile payments, peer-to-peer lending, digital wallets, and automated investment platforms. Fintechs are often designed for speed and convenience, which can create AML vulnerabilities. Many fintechs have rapid onboarding with minimal identity verification, making them attractive to criminals. Fintechs must comply with AML regulations and have been the target of regulatory enforcement actions. Under FATF, fintechs are treated as financial institutions or MSBs depending on their activities.

\vspace{0.1cm}
\noindent\textbf{Type of violation or warning signs:} A fintech onboards customers with only an email address and phone number (no identity verification). A fintech allows transfers to unverified recipients or to high-risk jurisdictions. A fintech has no transaction monitoring system. A fintech uses peer-to-peer networks that obscure the flow of funds. A fintech has no AML program and has never filed a SAR. A fintech is registered in a jurisdiction with weak AML controls but operates globally.

\vspace{0.1cm}
\noindent\textbf{Correct answer approach:} The correct answer will state that fintechs must comply with AML regulations including CDD, transaction monitoring, and SAR filing. Eliminate answers that suggest fintechs are exempt because they are "technology companies" not financial institutions. Remember: the activity determines the AML obligation, not the company's self-designation. If it transmits money, it's an MSB.
\end{frame}

% ============================================================
% SLIDE 5 - INSURANCE COMPANIES (LIFE AND GENERAL)
% ============================================================
\begin{frame}{5. Insurance Companies ML Risks}
\tiny
\noindent\textbf{Introduction to the topic:} Insurance companies, particularly life insurance and annuity providers, are vulnerable to money laundering through single-premium policies, policy loans, and early surrender. Criminals purchase a policy with illicit funds, take a policy loan to access clean funds, and then surrender the policy for additional clean funds. General insurance (property, casualty) can also be misused through fraudulent claims or over-insurance. Under FATF Recommendation 22, insurance companies must conduct CDD on policyholders, identify beneficial owners, and file SARs for suspicious activity. The amount of the premium and the source of funds are key risk indicators.

\vspace{0.1cm}
\noindent\textbf{Type of violation or warning signs:} A customer purchases a single-premium life insurance policy with \$10 million from an offshore account. A customer takes a policy loan shortly after purchase and wires the proceeds to a third party. A customer surrenders a policy early (within 6-12 months) after taking a loan. A customer makes multiple premium payments just below reporting thresholds. A customer's declared income cannot explain the premium amount. An insurance company does not verify the source of funds for large premiums.

\vspace{0.1cm}
\noindent\textbf{Correct answer approach:} The correct answer will state that insurance companies must conduct CDD and file SARs for suspicious activity, particularly for high-value policies, policy loans, and early surrenders. Eliminate answers that suggest insurance companies are exempt from AML rules. Remember: the pattern of purchase, loan, and early surrender is classic insurance money laundering.
\end{frame}

% ============================================================
% SLIDE 6 - SECURITIES FIRMS AND BROKER-DEALERS
% ============================================================
\begin{frame}{6. Securities Firms and Broker-Dealers}
\tiny
\noindent\textbf{Introduction to the topic:} Securities firms and broker-dealers facilitate trading in stocks, bonds, options, and other securities. They are vulnerable to money laundering through insider trading, market manipulation, and layering through rapid trading. Broker-dealers must conduct CDD on customers, monitor for suspicious trading patterns, and file SARs. Under SEC and FINRA rules, broker-dealers must have AML programs, designate compliance officers, and provide employee training. Cross-border brokerage accounts and omnibus accounts present additional risks because the broker may not know the underlying customer.

\vspace{0.1cm}
\noindent\textbf{Type of violation or warning signs:} A customer with no history of profitable trading suddenly has highly profitable trades before corporate announcements (insider trading). A customer executes trades that consistently lose money (layering). A customer trades through an omnibus account where the broker does not know the underlying customer. A customer moves funds between brokerage accounts in different countries without investment purpose. A customer uses wash trading to create artificial volume. A broker-dealer has no surveillance system for detecting suspicious trading.

\vspace{0.1cm}
\noindent\textbf{Correct answer approach:} The correct answer will state that broker-dealers must have AML programs, conduct CDD, monitor for suspicious trading, and file SARs. Eliminate answers that suggest broker-dealers have no AML obligations or that surveillance is optional. Remember: insider trading and market manipulation are predicate offenses for money laundering.
\end{frame}

% ============================================================
% SLIDE 7 - INVESTMENT ADVISERS AND FUND MANAGERS
% ============================================================
\begin{frame}{7. Investment Advisers and Fund Managers}
\tiny
\noindent\textbf{Introduction to the topic:} Investment advisers and fund managers manage assets for high-net-worth clients and institutional investors. They are vulnerable because they may not directly see the source of client funds (clients deposit funds directly into the manager's account). Advisers may also set up complex structures (hedge funds, private equity funds) that obscure beneficial ownership. Under FATF, investment advisers are subject to AML obligations including CDD on clients, identifying beneficial owners, monitoring for suspicious activity, and filing SARs. The source of wealth and source of funds for clients must be verified.

\vspace{0.1cm}
\noindent\textbf{Type of violation or warning signs:} An investment adviser accepts funds from a client without verifying source of wealth. A client invests through a shell company or trust with anonymous beneficiaries. A client's declared income cannot explain the investment amount. A client is a foreign PEP but the adviser does not conduct EDD. An adviser sets up a fund in a secrecy jurisdiction with no independent oversight. An adviser does not monitor withdrawals from the fund for suspicious activity.

\vspace{0.1cm}
\noindent\textbf{Correct answer approach:} The correct answer will state that investment advisers must conduct CDD, identify beneficial owners, and file SARs. Eliminate answers that suggest advisers have no AML obligations because they are not "banks." Remember: investment advisers are gatekeepers for the asset management industry — they must know the source of client funds.
\end{frame}

% ============================================================
% SLIDE 8 - CONSUMER LENDERS (PAYDAY, INSTALLMENT, MICROFINANCE)
% ============================================================
\begin{frame}{8. Consumer Lenders ML Risks}
\tiny
\noindent\textbf{Introduction to the topic:} Consumer lenders (payday lenders, installment lenders, microfinance institutions) provide small loans to individuals. While the loans themselves may be small, the cumulative volume can be large. Criminals may use consumer lenders to launder funds by taking out loans and then repaying with illicit funds (which are now "clean" repayments). Microfinance institutions in developing countries may have weak AML controls. Lenders must conduct CDD on borrowers, monitor for patterns of unusual activity, and file SARs for suspicious behavior.

\vspace{0.1cm}
\noindent\textbf{Type of violation or warning signs:} A borrower takes out multiple loans from different lenders and repays them early in cash. A borrower's declared income cannot explain the repayment amounts. A borrower requests that loan proceeds be sent to a third party rather than to themselves. A borrower uses a foreign passport or identification that cannot be verified. A microfinance institution does not verify borrower identity. A lender does not have a process for identifying suspicious repayment patterns.

\vspace{0.1cm}
\noindent\textbf{Correct answer approach:} The correct answer will state that consumer lenders must conduct CDD, monitor for patterns of unusual activity, and file SARs. Eliminate answers that suggest consumer lenders are exempt because loans are small. Remember: the cumulative volume matters — multiple small loans can be used for money laundering.
\end{frame}

% ============================================================
% SLIDE 9 - VIRTUAL ASSET SERVICE PROVIDERS (VASPs)
% ============================================================
\begin{frame}{9. Virtual Asset Service Providers (VASPs)}
\tiny
\noindent\textbf{Introduction to the topic:} Virtual Asset Service Providers (VASPs) exchange virtual assets for fiat, transfer virtual assets, or provide custody/administration of virtual assets. Under FATF, VASPs are subject to the same AML obligations as financial institutions. This includes registration, CDD, transaction monitoring, Travel Rule (sharing originator/beneficiary information), and SAR filing. VASPs are high-risk due to the pseudonymous nature of blockchain transactions and the use of privacy coins and mixers. Many VASPs have been fined for inadequate AML controls.

\vspace{0.1cm}
\noindent\textbf{Type of violation or warning signs:} A VASP has no registration or license. A VASP does not conduct KYC on customers (just email address). A VASP does not implement the Travel Rule. A VASP does not screen for sanctions or monitor for suspicious transactions. A VASP allows customers to use privacy coins (Monero) or mixers without enhanced due diligence. A VASP has never filed a SAR. A VASP is registered in a jurisdiction with weak AML controls.

\vspace{0.1cm}
\noindent\textbf{Correct answer approach:} The correct answer will state that VASPs must comply with AML regulations including registration, CDD, Travel Rule, and SAR filing. Eliminate answers that suggest VASPs are exempt because they are "decentralized" or not financial institutions. Remember: FATF guidance explicitly applies AML obligations to VASPs.
\end{frame}

% ============================================================
% SLIDE 10 - CRYPTOCURRENCY EXCHANGES
% ============================================================
\begin{frame}{10. Cryptocurrency Exchanges}
\tiny
\noindent\textbf{Introduction to the topic:} Cryptocurrency exchanges allow customers to buy, sell, and trade cryptocurrencies. They are a primary gateway between fiat and crypto and are therefore critical for AML compliance. Exchanges must conduct CDD on customers, monitor for suspicious transactions, screen for sanctions, and file SARs. High-risk activities include trading in privacy coins, use of mixers, rapid chain hopping, and transactions with unhosted wallets. Exchanges must also implement the Travel Rule to share information on transactions above the threshold.

\vspace{0.1cm}
\noindent\textbf{Type of violation or warning signs:} An exchange does not verify customer identity before allowing trading. An exchange allows large withdrawals to unhosted wallets without verification. An exchange does not screen for sanctions. An exchange does not monitor for suspicious patterns (rapid in-and-out, use of mixers). An exchange has no Travel Rule implementation. An exchange has never filed a SAR. An exchange uses a third-party provider for CDD without verifying the provider's controls.

\vspace{0.1cm}
\noindent\textbf{Correct answer approach:} The correct answer will state that cryptocurrency exchanges must conduct CDD, implement Travel Rule, monitor for suspicious activity, and file SARs. Eliminate answers that suggest exchanges are exempt from AML rules. Remember: exchanges are gatekeepers between fiat and crypto — they have the highest AML obligations in the crypto ecosystem.
\end{frame}

% ============================================================
% SLIDE 11 - PEER-TO-PEER LENDING PLATFORMS
% ============================================================
\begin{frame}{11. Peer-to-Peer Lending Platforms}
\tiny
\noindent\textbf{Introduction to the topic:} Peer-to-peer (P2P) lending platforms connect individual borrowers with individual lenders, bypassing traditional banks. While P2P lending can be legitimate, criminals may use these platforms to launder funds. For example, a criminal can lend illicit funds to a borrower, who then repays the loan (giving the criminal clean funds). Or a criminal can borrow funds and then invest them in other illicit activities. P2P platforms must conduct CDD on both lenders and borrowers, monitor for suspicious patterns, and file SARs. The source of funds for lenders and the use of loan proceeds by borrowers are key risk indicators.

\vspace{0.1cm}
\noindent\textbf{Type of violation or warning signs:} A lender's declared income cannot explain the volume of loans they are making. A borrower takes multiple loans from different platforms simultaneously. A lender funds loans with funds from a high-risk jurisdiction. A borrower repays a large loan early in cash. A P2P platform does not verify the identity of lenders or borrowers. A P2P platform has no transaction monitoring system. A P2P platform has never filed a SAR.

\vspace{0.1cm}
\noindent\textbf{Correct answer approach:} The correct answer will state that P2P platforms must conduct CDD, monitor for suspicious patterns, and file SARs. Eliminate answers that suggest P2P lending is an exception to AML rules. Remember: P2P platforms are financial intermediaries — they must know their customers and monitor for misconduct.
\end{frame}

% ============================================================
% SLIDE 12 - MONEY MARKET MUTUAL FUNDS
% ============================================================
\begin{frame}{12. Money Market Mutual Funds}
\tiny
\noindent\textbf{Introduction to the topic:} Money market mutual funds (MMFs) invest in short-term, low-risk securities and are often used by corporations and institutions as cash management vehicles. Criminals may use MMFs to park illicit funds temporarily — the funds are invested in a legitimate, low-risk vehicle and can be withdrawn as clean money later. MMFs must conduct CDD on investors, monitor for unusual patterns, and file SARs. Large inflows and rapid redemptions are particularly suspicious, especially when followed by wire transfers to high-risk jurisdictions.

\vspace{0.1cm}
\noindent\textbf{Type of violation or warning signs:} A customer deposits a large amount into an MMF, holds it for a short period, and redeems it. A customer deposits funds from an offshore account into an MMF and redeems to a different account. A customer's declared income cannot explain the deposit amount. A customer makes multiple deposits just below reporting thresholds. An MMF does not verify the source of funds for large deposits. An MMF does not monitor for rapid in-and-out activity.

\vspace{0.1cm}
\noindent\textbf{Correct answer approach:} The correct answer will state that MMFs must conduct CDD, monitor for suspicious activity, and file SARs. Eliminate answers that suggest MMFs are exempt from AML rules because they are "regulated" or "low-risk." Remember: any financial vehicle can be misused for money laundering — MMFs are no exception.
\end{frame}

% ============================================================
% SLIDE 13 - FACTORING AND TRADE FINANCE COMPANIES
% ============================================================
\begin{frame}{13. Factoring and Trade Finance Companies}
\tiny
\noindent\textbf{Introduction to the topic:} Factoring and trade finance companies provide financing based on accounts receivable or trade transactions. They are vulnerable to money laundering through over-invoicing (invoice is inflated to move funds), under-invoicing (goods are undervalued to move funds), and duplicate financing (same invoice is financed multiple times). Criminals use trade-based money laundering (TBML) to move funds across borders disguised as legitimate trade. Factoring companies must conduct due diligence on their clients and the underlying trade transactions, and file SARs for suspicious activity.

\vspace{0.1cm}
\noindent\textbf{Type of violation or warning signs:} A client's invoices are significantly above market value (over-invoicing). A client's invoices are from a high-risk jurisdiction. A client has no physical operations or inventory. A client uses multiple factoring companies to finance the same invoices. A client requests that proceeds be sent to an offshore account. A factoring company does not verify the legitimacy of the underlying trade transaction. A factoring company has no AML program.

\vspace{0.1cm}
\noindent\textbf{Correct answer approach:} The correct answer will state that factoring companies must conduct due diligence on clients and transactions, and file SARs for suspicious activity. Eliminate answers that suggest factoring companies are exempt from AML rules. Remember: trade-based money laundering is a common technique — factoring companies are on the front line.
\end{frame}

% ============================================================
% SLIDE 14 - NONBANK LENDING (MORTGAGE, COMMERCIAL, ASSET-BASED)
% ============================================================
\begin{frame}{14. Nonbank Lending (Mortgage, Commercial, Asset-Based)}
\tiny
\noindent\textbf{Introduction to the topic:} Nonbank lenders (mortgage companies, commercial lenders, asset-based lenders) provide loans without accepting deposits. They are vulnerable because they may have less regulatory oversight than banks, and they may not have robust AML programs. Criminals may use nonbank lenders to obtain loans with fraudulent collateral, or to repay loans with illicit funds (which are now "clean"). Lenders must conduct due diligence on borrowers, verify collateral, and file SARs for suspicious activity.

\vspace{0.1cm}
\noindent\textbf{Type of violation or warning signs:} A borrower's declared income cannot explain the loan amount. A borrower uses a shell company as collateral without disclosing beneficial ownership. A borrower repays a large loan early in cash. A borrower requests that loan proceeds be sent to a third party. A borrower has no legitimate business operations. A nonbank lender does not verify the source of funds for down payments or repayments. A nonbank lender has no AML program.

\vspace{0.1cm}
\noindent\textbf{Correct answer approach:} The correct answer will state that nonbank lenders must conduct CDD, verify the source of funds, and file SARs. Eliminate answers that suggest nonbank lenders are exempt from AML rules because they are not banks. Remember: any institution that extends credit is a potential money laundering gateway.
\end{frame}

% ============================================================
% SLIDE 15 - REGULATION AND OVERSIGHT OF NBFIs
% ============================================================
\begin{frame}{15. Regulation and Oversight of NBFIs}
\tiny
\noindent\textbf{Introduction to the topic:} NBFIs are regulated by various agencies depending on their activities. Money transmitters and MSBs are regulated by FinCEN and state regulators. Securities firms are regulated by SEC and FINRA. Insurance companies are regulated by state insurance departments. VASPs are regulated by FinCEN and may require state licenses. All NBFIs must register with the appropriate regulator, maintain AML programs, and file SARs. NBFIs are subject to examination and enforcement actions for AML deficiencies. Regulators look for registration, AML programs, CDD, transaction monitoring, and SAR filing.

\vspace{0.1cm}
\noindent\textbf{Type of violation or warning signs:} An NBFI has not registered with the appropriate regulator. An NBFI has no AML program, no compliance officer, and no training. An NBFI has never filed a SAR despite high volume. An NBFI has been cited for AML deficiencies in a previous examination but has not corrected them. An NBFI operates in a jurisdiction with weak AML controls. An NBFI's AML program is not commensurate with its risk profile.

\vspace{0.1cm}
\noindent\textbf{Correct answer approach:} The correct answer will state that NBFIs must register, maintain AML programs, and comply with regulatory oversight. Eliminate answers that suggest NBFIs are not subject to examination or enforcement. Remember: NBFIs are regulated — they must register and comply with AML requirements. Failure to do so results in enforcement action.
\end{frame}

% ============================================================
% SUMMARY SLIDE - NONBANK FINANCIAL AML PRINCIPLES
% ============================================================
\begin{frame}{Summary: Key Principles for Nonbank Financial ML Risks}
\tiny
\noindent\textbf{How to answer exam questions on nonbank financial services:}

\vspace{0.1cm}
\noindent• \textbf{MSBs/Money Transmitters:} Must register with FinCEN, maintain AML programs, and file CTRs and SARs.

\vspace{0.1cm}
\noindent• \textbf{Payment Processors:} Must conduct due diligence on merchants, monitor transactions, and file SARs.

\vspace{0.1cm}
\noindent• \textbf{Fintech:} Must comply with AML regulations — "technology company" is not an exemption from AML rules.

\vspace{0.1cm}
\noindent• \textbf{Insurance Companies:} Must conduct CDD on policyholders; pattern of purchase, loan, surrender = red flag.

\vspace{0.1cm}
\noindent• \textbf{Securities Firms:} Must have AML programs, surveillance systems, and file SARs for suspicious trading.

\vspace{0.1cm}
\noindent• \textbf{Investment Advisers:} Must verify source of wealth and source of funds for clients.

\vspace{0.1cm}
\noindent• \textbf{Consumer Lenders:} Must conduct CDD and monitor for suspicious repayment patterns.

\vspace{0.1cm}
\noindent• \textbf{VASPs/Crypto Exchanges:} Must register, implement Travel Rule, conduct CDD, monitor, and file SARs.

\vspace{0.1cm}
\noindent• \textbf{P2P Lending:} Must conduct CDD on both lenders and borrowers and monitor for suspicious patterns.

\vspace{0.1cm}
\noindent• \textbf{MMFs:} Must monitor for rapid in-and-out activity and file SARs for suspicious deposits/redemptions.

\vspace{0.1cm}
\noindent• \textbf{Factoring/Trade Finance:} Must verify legitimacy of trade transactions and file SARs for TBML.

\vspace{0.1cm}
\noindent• \textbf{Nonbank Lending:} Must verify source of funds and collateral, and file SARs.

\vspace{0.1cm}
\noindent• \textbf{Registration:} All NBFIs must register with appropriate regulators.

\vspace{0.1cm}
\noindent• \textbf{SAR Filing:} NBFIs must file SARs when they suspect money laundering.

\vspace{0.1cm}
\noindent• \textbf{Examination:} NBFIs are subject to regulatory examination and enforcement.

\vspace{0.3cm}
\centering
{Remember: Nonbank financial services are not "shadow banking" — they are regulated and must comply with AML rules. When in doubt, register, verify, monitor, and report.}
\end{frame}

\end{document}